% !TEX TS-program = xelatex
% !TEX encoding = UTF-8 Unicode
% !Mode:: "TeX:UTF-8"

\documentclass{resume}
\usepackage{zh_CN-Adobefonts_external}
\usepackage{linespacing_fix}
\usepackage{cite}

\begin{document}
\pagenumbering{gobble}

\name{侯汶政｜Go后端工程师}

\section{\faInfoCircle\ 基本信息}
\vspace{0.5ex} % 增加一点间距避免拥挤
\begin{center}
\begin{tabular}{p{0.45\textwidth} p{0.45\textwidth}}
    \textbf{个人网站：}\href{https://Soofjan.com}{Soofjan.com}\\
    
    \textbf{微信：}1489938120          & \textbf{当前状态：}深圳在职 \\
    \textbf{邮箱：}Soofjan1489938120@gmail.com & \textbf{到岗时间：}\textbf{2周} \\
    \textbf{学历：}本科全日制（\textbf{学信网可查}） & \textbf{求职偏好：}不外包 / 不驻场 \\
\end{tabular}
\end{center}

\section{\faCogs\ 技术能力}
\begin{itemize}[parsep=0.3ex]
    \item \textbf{编程语言与框架}：使用 Go（Gin、Gorm） 开发，系统掌握 关键字、并发、GMP 调度、内存管理等知识；拥有 Vue 实战经验，具备多端适应能力。
    \item \textbf{数据存储与缓存}：掌握 MySQL 存储引擎（InnoDB）、B+ 树索引、锁与事务（ACID、隔离级别、MVCC）；掌握 Redis 底层数据结构以及内部实现，熟悉过期策略、持久化、缓存等相关问题。
    \item \textbf{中间件}：熟悉 NSQ 架构，能处理消息丢失与重复消费；了解Kafka、RocketMQ等消息队列在不同场景下的使用选择。
    \item \textbf{开发工具}：熟悉 Git/Linux 开发环境；熟悉 TCP/IP、HTTP/HTTPS 等网络基础；高效使用 Cursor AI 工具。
\end{itemize}

\section{\faGraduationCap\ 教育背景}
\datedsubsection{\textbf{广州大学} \quad 学士，计算机科学与技术}{2019.09 -- 2023.06}
CET-6 英语六级（602分），全国大学生算法设计与编程挑战赛（银奖）

\section{\faSuitcase\ 工作经历}
\datedsubsection{\textbf{软件工程师 @ 慧为智能科技有限公司}}{2023.05 - 至今}

\paragraph*{私有云存储系统（NAS）开发}
% \mbox{} \\
% \textcolor{gray}{ \textbf{技术栈：} Go, Vue, C++, MySQL, SQLite, NSQ, AES+RSA, ElasticSearch, Docker  }

\begin{itemize}

    \item \textbf{云端}：
    \begin{itemize}
        \item 主导设计 NAT 穿透交互架构，规范客户端、服务端、云端及穿透端的通信流程，实现复杂网络下稳定高效访问。
        \item 引入 NSQ 消息队列，实现事件异步解耦，提升吞吐1.6倍，高并发平均延迟降低约 50\%。
        \item 采用应用层 AES+RSA 和传输层 TLS 的双层加密方案，保障凭证与敏感字段在复杂网络下的传输安全。
    \end{itemize}

    \item \textbf{智能文档中心}：
    \begin{itemize}
        \item 针对内存受限环境，完成了从 Elasticsearch 到 Bleve 的搜索引擎迁移，结合 GPT 实现快速的全文与多维度检索。
        \item 设计基于块的索引策略与搜索结果聚合算法，解决超长文档索引阶段内存占用过高问题，峰值 heapInuse 降幅约 76\%+。
    \end{itemize}
    
    \item \textbf{文件索引同步系统}：
    \begin{itemize}
        \item 自研 BFS+Queue 架构替代 fsnotify 监听，在海量文件场景下内存开销降低 90\% 以上。
        \item 基于接口抽象与 Functional Options 模式，实现同步逻辑与底层存储、IO 的深度解耦，支持跨项目的高内聚集成。
    \end{itemize}
    
    \item \textbf{其他}：
    \begin{itemize}
        \item AI相册: 在RKNN设备上实现物品分类；通过 DBSCAN 聚类算法实现万级人脸库的无监督聚类，支持动态聚类中心更新，识别准确率达 97.3\%。
        \item 通过 rate 限流、Ristretto 缓存库的方式有效缓解并发流量冲击，核心接口缓存命中率达 85\%，响应时间缩短 90ms。
        \item 开发并维护 Docker、虚拟机、Webdav 等前后端项目，提供 NAS 生态支持。
    \end{itemize}

\end{itemize}

\end{document}